%%%%%%%%%%%%%%%%%%%%%%%%%%%%%%%%%%%%%%%%%%%%%%%%%%%%%%%%%%%%%%%%%%%%%%%%%%%%%%%%
%
%   大连理工大学本科毕业设计（论文）模板 —— 封面文件 cover.tex
%   适配版本：1.0（基于原版0.71优化）
%   适配场景：大连理工大学本科毕业设计（论文）封面及摘要页
%   编译环境：CTeXLive 2021+ / Overleaf / Windows/macOS/Linux
%
%%%%%%%%%%%%%%%%%%%%%%%%%%%%%%%%%%%%%%%%%%%%%%%%%%%%%%%%%%%%%%%%%%%%%%%%%%%%%%%%

% 核心封面信息配置（按需修改，字符间距已优化适配封面排版）
\cdegree{大连理工大学本科毕业设计（论文）}
\ctitle{基于\texttt{COMTool}项目的代码质量与安全分析} % 示例标题，替换为你的论文标题
\etitle{Code Quality and Security Analysis Based on COMTool Project} % 英文标题

% 学院/专业/作者信息（空格为封面排版对齐用，可根据实际姓名长度微调）
\cschool{能~源~与~动~力~学~院}
\csubject{计~算~数~学~专~业}
\cauthor{张~~三}
\cauthorno{200801000}
\csupervisor{王~五　副教授}
\creviewer{张~三　教授}

% 日期：默认使用编译时间，也可手动指定（例：\cdate{2024~年~06~月~10~日}）
\cdate{\the\year~年~\the\month~月~\the\day~日}

% 中文摘要（贴合本科毕设要求：600-800字，含目的、方法、结果、结论，无插图）
\cabstract{
本设计以开源项目COMTool为研究对象，围绕代码质量与软件安全展开系统性分析。首先基于radon库完成代码复杂度评估，量化圈复杂度、可维护性指数等核心指标；其次通过pipdeptree与safety工具开展依赖安全审计，排查供应链安全风险；再利用bandit静态扫描工具识别代码层面的安全漏洞，并按风险等级分类；最后结合PyGithub库分析GitHub生态数据，涵盖Issues、PR及Actions工作流等维度。

分析结果表明：COMTool项目整体代码复杂度控制良好，无已知漏洞依赖，但存在5个高风险代码安全漏洞及1个高风险GitHub安全Issue。基于分析结论，提出针对性的优化建议，为开源项目的代码质量提升与安全加固提供实践参考。本设计验证了自动化代码分析工具在开源项目评估中的有效性，也为同类软件的质量管控提供了可复用的分析方法。

关键词选取遵循《汉语主题词表》规范，聚焦研究核心内容，共选取3-5个核心关键词。
}

% 中文关键词（3-5个，分号分隔，末尾无标点）
\ckeywords{代码质量分析；软件安全；静态扫描；COMTool；开源项目}

% 英文摘要（与中文摘要对应，简洁准确）
\eabstract{
This design takes the open-source project COMTool as the research object, and carries out a systematic analysis around code quality and software security. Firstly, the code complexity evaluation is completed based on the radon library to quantify core indicators such as cyclomatic complexity and maintainability index. Secondly, dependency security audit is carried out through pipdeptree and safety tools to investigate supply chain security risks. Then, the bandit static scanning tool is used to identify security vulnerabilities at the code level and classify them by risk level. Finally, GitHub ecological data including Issues, PR and Actions workflows are analyzed based on the PyGithub library.

The analysis results show that the overall code complexity of the COMTool project is well controlled, with no known vulnerable dependencies, but there are 5 high-risk code security vulnerabilities and 1 high-risk GitHub security Issue. Based on the analysis conclusions, targeted optimization suggestions are put forward to provide practical reference for the code quality improvement and security reinforcement of open-source projects. This design verifies the effectiveness of automated code analysis tools in the evaluation of open-source projects, and also provides reusable analysis methods for quality control of similar software.
}

% 英文关键词（与中文关键词对应，逗号分隔）
\ekeywords{Code Quality Analysis, Software Security, Static Scanning, COMTool, Open-source Project}

% 生成封面（请勿修改此行）
\makecover